\section{Related Work}
\label{sec:related_work}

\paragraph{Classical training-free anomaly detection.}
A long line of work builds anomaly detectors on top of frozen pretrained representations, scoring query patches by distance to a memory bank of normal features. PatchCore~\citep{roth2022patchcore} established nearest-neighbour scoring on ImageNet-pretrained CNN features; later work replaced the backbone with DINOv2~\citep{oquab2024dinov2} and demonstrated stronger cross-domain transfer~\citep{subspacead2026}. Zero-shot approaches such as WinCLIP~\citep{jeong2023winclip} and AnomalyCLIP~\citep{zhou2024anomalyclip} cast the problem as similarity between image regions and textual prototypes (``normal'' vs.\ ``defective''). These methods share a common assumption: the anomaly signal is visible in the embedding space of a generic pretrained feature extractor, without any reasoning step. AnomalyClaw builds directly on this line: SubspaceAD and a DINOv2 patch-$k$-NN provide the expert evidence that the adaptive router consumes. Unlike classical detectors, the expert never makes the final decision by itself on VLM-uncertain items; instead, its score statistics ($\rho$, $\kappa$) drive routing, and its flagged patches become text evidence for the VLM's re-examination call.

\paragraph{VLM-based anomaly detection.}
Recent work uses large VLMs directly for anomaly decisions. AnomalyGPT~\citep{gu2024anomalygpt} fine-tunes a VLM on industrial defect data with simulated anomalies; MMAD~\citep{jiang2024mmad} evaluates several frontier VLMs on a multi-task anomaly benchmark. The dominant deployment pattern, however, remains a \emph{single forward pass} with a short prompt asking whether the image is abnormal, optionally with one or two reference images. As we show in \S\ref{sec:experiments}, this pattern is substantially under-calibrated when paired with generic domain descriptions: a one-sentence task-anchored descriptor alone moves a frontier VLM from $0.761$ to $0.825$ macro AUROC. AnomalyClaw stays on top of such a VLM without any fine-tuning and treats the descriptor-enhanced single-pass call as its first phase.

\paragraph{Agentic anomaly detection.}
A recent wave of work adds multi-turn reasoning, tool use, retrieval, and memory on top of a base VLM. AgentIAD~\citep{zhang2025agentiad} orchestrates a VLM with a patch expert via a tool interface; EAGLE~\citep{chen2026eagle} and similar systems introduce multi-role debate between VLM-based agents; AutoIAD~\citep{wei2025autoiad} adds a routing controller that selects methods per image. These systems report gains over single-pass baselines, but three structural issues make the gains hard to attribute: (a) they bundle \emph{what to look at} (tools), \emph{who provides evidence} (expert choice), and \emph{how to aggregate} (strategy) into a single monolithic pipeline; (b) they typically compare to a generic-descriptor VLM baseline that we show is badly under-calibrated; (c) extra VLM calls, retrieval indices, per-domain calibration, and multi-turn exchange all confound the ablation. AnomalyClaw makes three deliberate choices. First, the framework \emph{decouples} the three decisions (Tool, Expert, Strategy) so that each can be swapped independently and reported separately. Second, our baselines include both the descriptor-enhanced single-pass call (a strong baseline most prior agentic work did not report) and a zero-cost score-fusion baseline that already matches or exceeds every single agent strategy on our benchmark. Third, our router picks one of four strategies via calibration-argmax or descriptor rules---it does not invent new inference on top of the chosen strategy---and the ablation in Appendix~\ref{app:ablation} documents six alternative agent designs that all \emph{hurt} the stronger baseline on GPT-5.4 calibration.

\paragraph{Self-refinement and debate.}
Outside VAD, self-refinement~\citep{madaan2023selfrefine} and debate~\citep{du2023debate,irving2018debate} have been proposed as general recipes for improving LLM reliability. Applied naively on top of a strong descriptor-enhanced VLM baseline, however, both recipes regress rather than improve: a second pass prompted to question the first rationalises correct high-confidence predictions away. Our self-refine and debate ablations (\S\ref{ssec:ablation}) are included precisely to show that \emph{unconditional} re-examination is the wrong move. AnomalyClaw instead uses the expert signal as an \emph{asymmetric, evidence-driven} trigger: a second VLM call fires only when the expert flags a strong, localised deviation that the VLM missed.

\paragraph{Cross-domain VAD benchmarks.}
MVTec-AD~\citep{bergmann2019mvtec} and VisA~\citep{zou2022visa} remain the dominant benchmarks for training-free VAD, but both are industrial only. Remote-sensing~\citep{chen2020levircd}, medical~\citep{bao2024bmad,yang2023medmnist}, and road-anomaly~\citep{lis2019detecting,yu2020bdd100k} benchmarks exist in isolation, making it hard to argue about cross-domain generalisation. CrossDomainVAD-12 unifies twelve domains under a single reference-based image-level evaluation protocol with documented sampling, splits, and taxonomy, and is explicitly designed to surface the descriptor- and expert-complementarity questions that industrial-only evaluation hides.

\paragraph{Tool-using language agents.}
Outside VAD, ReAct~\citep{yao2023react} and Toolformer~\citep{schick2023toolformer} demonstrated that LLMs can be prompted to invoke an external tool registry. Our framework inherits the Tool-$\times$-Strategy decomposition from this tradition; the novelty is (a) applying it to training-free VAD with non-parametric experts as a first-class axis, and (b) showing that a \emph{lightweight, calibration-only} router outperforms free-form VLM tool-calling on our benchmark.

\paragraph{Positioning.}
Our contribution is (i) a decoupled Tool-$\times$-Expert-$\times$-Strategy framework for cross-domain VAD, (ii) a calibration-only router that outperforms every single-strategy baseline on all three backbones (significantly on SeedVL), (iii) an empirical study of what moves training-free VAD performance on top of modern VLMs, and (iv) a benchmark and a documented protocol to make such claims auditable.
